% Direct statement of the Jacobian Challenge.
%
% The blocks here are the informal counterparts of the public Lean
% obligations in `Jacobian/Challenge.lean`. They fix the genus, the
% Jacobian, the Abel-Jacobi map, the degree, pushforward/pullback, and
% the four anti-hack theorems that the formalization must discharge.
% Internal helper statements (period pairing, period lattice, trace on
% one-forms, etc.) live in the Statement Bank that follows.

\section{Challenge Statement}
\label{sec:challenge-statement}

Throughout, \(X\) is a compact connected Riemann surface, i.e.\ a
nonempty compact connected Hausdorff space carrying an analytic complex
manifold atlas (\code{ChartedSpace ℂ X} together with
\code{IsManifold 𝓘(ℂ) $\omega$ X} in the Lean target). The challenge fixes
six pieces of data and four anti-hack theorems on top of that setup.

\subsection{Genus}

\begin{definition}[Analytic genus]
\label{def:analytic-genus}
For a compact connected Riemann surface \(X\), define
\[
  g(X) := \dim_{\C} H^{0}(X,\Omega^{1}_{X}).
\]
\lean{JacobianChallenge.HolomorphicForms.analyticGenus}
\leanok
\proofplan
Instantiate finite-dimensionality of \(H^0(X,\Omega_X^1)\) from the Sec02 chain:
sup norm \(\to\) compact unit ball (Montel) \(\to\) Riesz theorem; then define
\(g(X)\) by `FiniteDimensional.finrank`.
\end{definition}

\begin{layman}
The \emph{analytic genus} of a compact Riemann surface is the complex
dimension of the space of global holomorphic 1-forms on it. The
finite-dimensionality theorem (proved later, in the holomorphic-forms
section) makes this an natural number. By the genus-zero classification below, this number also matches
the topological genus --- the count of handles you'd see if you held the
surface in your hands. The sphere has analytic genus 0, the donut has 1,
the two-handled pretzel has 2, and so on. Why bother? Genus is the
single most important invariant in Riemann surface theory. The whole
challenge revolves around tying the analytic definition (vector-space
dimension) to the topological one (number of handles), and the bridge
between them is forged by Riemann--Roch.
\end{layman}
\begin{theorem}[Genus-zero classification]
\label{thm:genus-zero-classification}
\uses{def:analytic-genus, thm:analytic-eq-topological-genus, lem:holomorphic-one-form-linearEquiv-to-OnePointCx, thm:homeo-sphere-holomorphic-one-form-vanishing, thm:degree-one-bijectivity-meromorphic}
For a compact connected Riemann surface \(X\),
\[
  g(X)=0 \quad\Longleftrightarrow\quad X \cong_{\mathrm{homeo}} \Sphere .
\]
\lean{genus_eq_zero_iff_homeo}
\leanok
\proofplan
\(\Rightarrow\): use Riemann--Roch at genus \(0\) to build degree-one meromorphic map
to \(\mathbf{P}^1\), then classify degree-one maps as homeomorphisms.
\(\Leftarrow\): transport the vanishing of holomorphic forms on \(\mathbf{P}^1\)
to \(X\) via the biholomorphism (Theorem~\ref{lem:holomorphic-one-form-linearEquiv-to-OnePointCx}).
\end{theorem}

\begin{layman}
This is the first anti-hack theorem of the challenge: the analytic
genus of a compact Riemann surface is zero if and only if the surface
is homeomorphic to the topological sphere. Forward direction: genus
zero produces a degree-one meromorphic function, which is a
biholomorphism onto the Riemann sphere, which is homeomorphic to the
2-sphere. Reverse direction: a homeomorphism with the sphere makes
the topological genus zero, and the comparison theorem identifies
analytic and topological genus, giving analytic genus zero. Why
bother? It forbids the trivial fake construction
``\code{genus X := 0}.'' The genus has to actually equal the
topological genus --- ruling out any constant function pretending to be
genus. Without this anchor, all the challenge's anti-hacks would be
wide open.
\end{layman}

\subsection{The Jacobian}

\begin{definition}[Analytic Jacobian]
\label{def:analytic-jacobian}
\uses{thm:period-lattice, prop:complex-torus-package, def:analytic-genus}
The Jacobian of \(X\) is the complex torus
\[
  \Jac(X) := \HdualX / \periods_X ,
\]
where \(\periods_X\) is the period lattice of Theorem~\ref{thm:period-lattice}.
The challenge requires \(\Jac(X)\) to carry, simultaneously, the
structures of an additive commutative group, a Hausdorff compact
topological space, a complex manifold of complex dimension \(g(X)\), and
a complex Lie additive group on that manifold structure.
\lean{Jacobian}
\leanok
\proofplan
Use period lattice theorem + complex torus quotient package to construct
\(\HdualX/\periods_X\) with additive, topological, and manifold structures, and
to install all the \code{AddCommGroup}/\code{TopologicalSpace}/\code{T2Space}/\code{CompactSpace}/\code{ChartedSpace}/\code{IsManifold}/\code{LieAddGroup}
instances required by \code{Jacobian/Challenge.lean}.
\end{definition}

\begin{layman}
The \emph{Jacobian} of a compact Riemann surface is the quotient: take
the dual of the holomorphic 1-forms (a finite-dimensional complex
vector space), mod out by the period lattice (a full discrete grid),
and the result is a compact complex torus --- a higher-dimensional
generalization of an ordinary donut. Why bother? This is the central
object of the challenge. Every theorem in the public API is about the
Jacobian or about maps into and out of it. The construction inherits
all the structural goodness of the complex-torus package --- addition,
compactness, holomorphic atlas, Lie group structure --- for free, just
by plugging in the dual space and the period lattice. The compact
complex Lie group structure on \(\Jac(X)\) is itself an anti-hack: any
construction that fails to deliver every layer (group, topology,
manifold, Lie) is rejected.
\end{layman}

\subsection{The Abel--Jacobi map}

\begin{definition}[Abel-Jacobi map]
\label{def:abel-jacobi}
\uses{def:analytic-jacobian, lem:aj-path-independence}
For \(P\in X\), define
\[
  \AJ_P(Q) :=
    \left[ \omega \mapsto \int_{P}^{Q}\omega \right] \in \Jac(X),
\]
where the path is arbitrary and the class is taken modulo periods. The
challenge requires \(\AJ_P\) to be holomorphic and to send \(P\) to the
identity element of \(\Jac(X)\).
\lean{Jacobian.ofCurve}
\leanok
\proofplan
Define by integrating forms from \(P\) to \(Q\); prove path-independence modulo
period lattice, then descend to Jacobian quotient class. Holomorphicity is
\code{Jacobian.ofCurve_contMDiff}; the identity-at-base-point obligation is
\code{Jacobian.ofCurve_self}.
\end{definition}

\begin{layman}
Pick a base point on the surface. The \emph{Abel--Jacobi} map sends
any other point to the class in the Jacobian represented by ``integrate
every holomorphic 1-form along a path from base to here.'' Different
paths give answers differing by periods, but the period lattice is
exactly what the Jacobian quotients out, so the class doesn't depend
on which path you chose. The result is a holomorphic map from the
surface into its own Jacobian, sending the base point to zero. Why
bother? This is the central comparison map between the surface and
its Jacobian. It's holomorphic, and (crucially, on positive-genus
surfaces) it's injective --- so the surface sits inside its Jacobian as a
curve rather than collapsing to a point.
\end{layman}

\begin{theorem}[Abel-Jacobi theorem, point-separation form]
\label{thm:abel-jacobi-injective}
\uses{def:abel-jacobi, input:abel-theorem}
If \(g(X)>0\), then \(\AJ_P : X\to \Jac(X)\) is injective.
\lean{Jacobian.ofCurve_inj}
\leanok
\proofplan
Reduce equality of Jacobian classes to vanishing of all period functionals and
invoke Abel point-separation (Abel theorem input).
\end{theorem}

\begin{layman}
On a compact Riemann surface of positive genus, the Abel--Jacobi map is
\emph{injective}: distinct points on the surface go to distinct points
in the Jacobian. This is the second anti-hack theorem of the
challenge. The proof chain (worked out in detail in the Abel--Jacobi
section):
equality of Abel--Jacobi images forces a principal divisor of the form
``one point minus another,'' which gives a degree-one map to the
Riemann sphere, which forces the surface to be the sphere --- but the
sphere has genus zero, contradiction. Why bother? Injectivity rules
out trivial fake Jacobians like ``\code{Jacobian X := PUnit}''
(which would collapse all points to one) or any other construction
that fails to actually distinguish the surface's points. The compact-complex-Lie-group structure on the Jacobian, plus
point-separation, forces the construction to separate points.
\end{layman}

\subsection{Degree, pushforward, and pullback}

\begin{definition}[Degree and trace]
\label{def:degree-trace}
\uses{lem:pullback-forms, lem:trace-forms, def:branched-degree, thm:branched-degree-well-defined}
For a holomorphic map \(f:X\to Y\), define \(\deg(f)\) to be \(0\) if
\(f\) is constant and otherwise the topological degree of the associated
branched covering. For nonconstant \(f\), define the trace on one-forms
\[
  \tr_f : H^{0}(X,\Omega^{1}_{X}) \to H^{0}(Y,\Omega^{1}_{Y})
\]
by summing local inverse branches away from the branch locus and extending
holomorphically. Only \(\deg(f)\) is exposed in the public API
(\code{ContMDiff.degree}); the trace is internal infrastructure used to
build the pushforward map below.
\lean{ContMDiff.degree, JacobianChallenge.TraceDegree.pullbackFormsLinearMap}
\leanok
\proofplan
Define degree by constant/nonconstant split; define trace by summing inverse
branches off branch locus and extending holomorphically across it.
\end{definition}

\begin{layman}
Two operations associated to a holomorphic map between compact Riemann
surfaces. The \emph{degree} counts how many times the map covers its
target --- zero if the map is constant; otherwise the size of a generic
fiber, with multiplicities at branch points. The \emph{trace on
1-forms} is the natural way to push a holomorphic 1-form on the source
down to a holomorphic 1-form on the target, by summing over local
inverse branches and extending holomorphically across the branch
locus. Both are core ingredients of the push-pull machinery the
challenge demands. Why bother? Degree is the ``size'' of a
holomorphic map; trace is the natural pushforward operation on
1-forms. Together with pullback, they assemble into the
pushforward-pullback identity below, which forces the right
interaction between push, pull, and degree on Jacobians.
\end{layman}

\begin{theorem}[Functoriality and holomorphicity of induced Jacobian maps]
\label{thm:push-pull-functoriality}
\uses{lem:push-pull-descend}
For every holomorphic map \(f : X \to Y\) of compact connected Riemann
surfaces, the induced maps
\[
  f_{*} : \Jac(X) \to \Jac(Y),\qquad
  f^{*} : \Jac(Y) \to \Jac(X)
\]
are holomorphic continuous group homomorphisms. They satisfy the
identity and composition laws
\[
  (\mathrm{id})_*=\mathrm{id},\qquad (g\circ f)_*=g_*\circ f_*,
\]
\[
  (\mathrm{id})^{*}=\mathrm{id},\qquad (g\circ f)^{*}=f^{*}\circ g^{*}.
\]
\lean{Jacobian.pushforward, Jacobian.pullback, Jacobian.pushforward_contMDiff, Jacobian.pushforward_id_apply, Jacobian.pushforward_comp_apply, Jacobian.pullback_contMDiff, Jacobian.pullback_id_apply, Jacobian.pullback_comp_apply}
\leanok
\proofplan
Construct descended maps from linear maps on \(\HdualX\), prove they preserve
period lattice, then establish holomorphicity and identity/composition laws.
\end{theorem}

\begin{layman}
The maps induced on Jacobians by holomorphic maps between Riemann
surfaces are well-behaved: they are holomorphic group homomorphisms,
they take the identity map of a surface to the identity map of its
Jacobian, and they compose correctly under composition of holomorphic
maps. The pullback reverses arrow direction (target to source); the
pushforward respects arrow direction (source to target). It's
``functoriality'' in the categorical sense: the assignment
``Riemann surface'' \(\mapsto\) ``Jacobian'' is a functor --- both
contravariantly (via pullback) and covariantly (via pushforward).
Why bother? Without functoriality, the push-pull operations would be
ad-hoc and could be gamed; with it, they're forced to behave like
geometric pushforwards and pullbacks. The challenge's API
requires this exact set of compatibility laws.
\end{layman}

\begin{theorem}[Pushforward-pullback identity]
\label{thm:pushforward-pullback}
\uses{thm:trace-pullback, lem:push-pull-descend}
For every \(P\in\Jac(Y)\),
\[
  f_{*}(f^{*}P)=\deg(f)\,P.
\]
\lean{Jacobian.pushforward_pullback}
\leanok
\proofplan
Descend the trace-pullback identity from one-forms to Jacobian quotient classes
and transport through the pushforward/pullback descent maps.
\end{theorem}

\begin{layman}
The fourth anti-hack theorem of the challenge. For any holomorphic
map between compact Riemann surfaces and any element of the target
Jacobian, applying pullback (to get an element of the source
Jacobian) and then pushforward (to get back to the target Jacobian)
recovers the original element multiplied by the degree of the map.
Combined with functoriality from the previous theorem, this forces
push, pull, and degree to interact through the classical trace
identity, no other interaction allowed. Why bother? This rules out
any fake construction that defines push and pull as arbitrary
unrelated homomorphisms. The geometric pullback and
pushforward must satisfy this identity, so any implementation in Lean
has to actually compute the trace correctly. Together with the genus
classification, point separation, and torus structure, these four
anti-hack theorems make trivial fake Jacobians impossible.
\end{layman}
